\ssr{ВВЕДЕНИЕ}

Существуют программы, в которых требуется обрабатывать большие объёмы данных или производить множество операций, зависимых от процессора. Пример такой программы --- визуализация сцены с использованием алгоритма трассировки лучей, в котором обрабатываются несколько десятков миллионов операций. Для того чтобы ускорить время выполнения программы, используется \textbf{параллельное исполнение алгоритмов}.

\textbf{Параллельное исполнение} --- способ, при котором происходит одновременное исполнение алгоритма в нескольких потоках.

Целью данной работы является разработать и провести сравнительный анализ последовательного и параллельного алгоритма трассировки лучей, используемого для генерации изображения. Для достижения цели поставлены следующие задачи:
\begin{itemize}
	\item описать последовательный алгоритм решения задачи;
	\item разработать параллельную версию алгоритма;
	\item реализовать обе версии алгоритма;
	\item выполнить сравнительный анализ зависимостей времени решения задач от размерности входа для реализации последовательного алгоритма и для реализации модифицированного алгоритма;
	\item выполнить сравнительный анализ зависимостей времени решения задач от размерности входа для реализации модифицированного алгоритма при $k \in \{1, 2, 4, 8, \dots, 8 \cdot q\}$, где $q$ --- количество логических ядер процессора;
	\item сформулировать рекомендацию о выборе k для решения задачи на ЭВМ выбранной автором работы архитектуры.
\end{itemize}
\clearpage
